\chapter{Reinforcement-learning harness variations}
\label{ch:harnessvar}

The harness is held constant across architectures, but several of its knobs are
themselves objects of study, because the deceptive $0.5$ floor of the seven-step task makes
the optimisation as much a part of the result as the network.

\section{Discount and horizon}
The discount $\gamma$ is matched to the horizon: $\gamma=0.95$ for the seven-step task and
$\gamma=0.98$ for the twenty-nine-step task, where credit must travel further from the
change to the cue. These are not finely tuned; the qualitative behaviour is robust across a
band, but a higher $\gamma$ on the longer task does help the cue's effect survive the delay.

\section{Critic particles}
The number of distributional ``particles'' (quantiles, or categorical atoms in the paper's
form) was varied from the published $15$ down to $5$. Fewer particles is a coarser value
distribution; on this low-bandwidth task the value differences the critic must resolve are
small (the gap between ``press now'' and ``wait'' near threshold), so particle count is a
genuine knob rather than a free choice. We default to $5$ for the short task and have run
$15$ on the long task; a systematic sweep is an open item.

\section{Target network: EMA versus hard copy}
The original working checkpoints used a periodic hard-copy target. The reproduction adds an
exponential-moving-average target ($d=0.995$). Whether the EMA helps, hurts, or is neutral
on this task is genuinely open: the working multiplicative$+$VAE model used EMA, but so
would a hard-copy version, and we have not isolated the difference. Because the lab's
canonical working models used the hard-copy target, an EMA-versus-hard-copy A/B on the
frozen-VAE base is on the near-term list --- it is a candidate, alongside the attention form
and head count, for any residual difference between our reproduction and the lab's other
networks.

\section{Exploration and the two collapses}
The recurring optimisation failure is collapse to a $0.5$ optimum, in one of two faces:
always-wait (never press) or always-press-at-the-deadline (press every trial at the change
frame). Both have policy entropy at zero and zero policy gradient --- the actor is frozen,
even while the critic keeps learning the $0.5$ value perfectly. The levers that matter are
maintaining exploration (entropy floor or coefficient, the burn-in length, the initial
action bias) and, decisively, fixing perception (Chapter~\ref{ch:perception}): once the
front-end can see the change, the discriminating policy is reachable and the collapse is
escaped. The lesson for the harness is that an apparently dead run is usually a perception
or exploration problem, not a critic problem.

\section{A resume footgun, fixed}
A practical note that cost real training time: the trainer's \code{resume} mode silently
no-ops if no explicit checkpoint path is set, starting fresh and discarding progress. This
was hardened in both trainers to auto-discover the latest checkpoint in the run directory
and to error loudly when \code{resume} finds nothing, rather than silently restarting.
Checkpoints live outside the synced repository to avoid mid-run corruption.
